\chapter{Conditional computation and routing}
\label{ch:routing}

Six tokens choose among three feed-forward networks. Each network has two available slots. Six tokens, six slots: should every token fit?

No. If their choices are $(0,0,0,2,1,1)$, expert 0 receives three requests while expert 2 receives one. A spare slot at expert 2 does not help the third token requesting expert 0 unless the routing policy permits reassignment. A capacity budget is not a guarantee of balanced demand.

Chapter 9 assembled a dense Transformer block. Every valid token used the same feed-forward parameters. Here we replace only that branch with conditional expert execution, then make selection, admission, weighting, gradients, and state ownership explicit. Our small reference is not a Switch Transformer training reproduction or an optimized distributed runtime. Its purpose is to reveal the contracts a future Evolutor implementation must preserve.

\section{Separate scoring, selection, admission, and weighting}
\index{mixture of experts}\index{router}\index{top-1 selection}
Let $Z\in\mathbb R^{T\times D}$ contain the branch inputs for $T$ token positions.
There are $E$ experts, each mapping $\mathbb R^D$ to $\mathbb R^D$. A router
$R\in\mathbb R^{D\times E}$ produces logits $H=ZR$. For token $t$,
\begin{equation}
p_{tj}=\frac{\exp(h_{tj})}{\sum_{r=0}^{E-1}\exp(h_{tr})},
\qquad k_t=\arg\max_j p_{tj}. \label{eq:route}
\end{equation}
Expert indices start at zero. Probabilities sum to one over experts, not tokens.
Softmax is evaluated by the library's numerically stable operation. A row-wise
constant shift in logits leaves the probabilities unchanged.

Top-1 selection chooses one expert index. It does not yet grant a capacity slot.
Let $v_t$ indicate a valid token and $a_t$ indicate admission. Our policy visits
tokens in their input order, excludes invalid positions, and admits a token if
its selected expert has fewer than $C$ occupied slots. It never reroutes overflow.
The branch result is
\begin{equation}
b_t=a_t p_{t k_t} E_{k_t}(z_t),\qquad y_t=u_t+b_t. \label{eq:branch}
\end{equation}
$u_t$ is the residual representation being updated. A dropped branch leaves
$y_t=u_t$; it does not remove the token from the sequence or from a later loss.

This separation follows the useful top-1, probability-weighted expert mechanism
described by Fedus, Zoph, and Shazeer \cite{switch}. Our deterministic
first-arrival implementation and test cases are deliberately small. We do not
inherit that paper's training, hardware, or quality results.

\Plate{01-routing}{Six fully specified probability rows. The selected experts request $(3,2,1)$ slots; capacity two gives admitted counts $(2,2,1)$. Token 2 is rejected after tokens 0 and 1 occupy expert 0. The vacant slot at expert 2 remains vacant under this no-rerouting policy.}{fig:routing}
\begin{center}\input{\Generated/trace-table.tex}\end{center}

\subsection{Ties and padding are part of the definition}
If logits tie, a result must still be chosen. This reference uses the first
maximal expert index, matching the documented \texttt{torch.argmax} convention
\cite{argmax}. A tie is not broken by a random coin or by the least loaded
expert. Those would be different policies with different reproducibility
requirements.

An invalid padding position may have logits and a selected index in the
intermediate tensors, but it consumes no slot, contributes no expert branch,
and does not enter the load-balancing denominator. Attention validity and
next-token target validity remain the separate Chapter 9 contracts.

The implementation returns both \emph{demand} before admission and occupied
\emph{used} slots after admission. In a continued cohort, used counts include
earlier chunks, while demand reports only the current call.
\Listing{plan}

The Python loop exposes control flow but synchronizes scalar decisions with the
host. It is intended for transparent CPU tests, not GPU throughput. Real
dispatch kernels can implement the same policy through prefix counts and
indexed communication. They must first agree on exactly which policy they are
implementing.

\section{Gather only admitted rows, then restore token identity}
\index{dispatch}\index{gather and scatter}\index{expert parameters}
Each expert is a feed-forward network:
\begin{equation}
E_j(z)=\operatorname{GELU}(zW_{1j}+b_{1j})W_{2j}+b_{2j}.
\end{equation}
$W_1$ has stacked shape $[E,D,F]$, $b_1$ has $[E,F]$, $W_2$ has $[E,F,D]$,
and $b_2$ has $[E,D]$, where $F$ is hidden width. These are independent
parameter sets. The word expert describes ownership and conditional use; it
does not establish meaningful specialization before training and evaluation.
\Listing{expert}

For expert $j$, form the ordered row index set
\begin{equation}
I_j=\{t:\ a_t=1,\ k_t=j\}.
\end{equation}
Gather $Z[I_j]$ with shape $[|I_j|,D]$, evaluate only that expert on those rows,
multiply row $t$ by $p_{tj}$, and scatter into the original token positions.
Top-1 admission makes the $I_j$ disjoint, so no token receives two writes.
A top-$k$ extension would require accumulation, not blindly reusing this
single-write assumption.

\Plate{02-dispatch}{From token-major inputs to expert minibatches and back. Expert 0 receives rows 0 and 1; expert 1 receives rows 4 and 5; expert 2 receives row 3. The output preserves original token order, including a zero branch at row 2. The formulas use the probabilities in Figure \ref{fig:routing}.}{fig:dispatch}
\Listing{dispatch}

The zero initialization retains a zero-valued autograd connection to inputs and
router probabilities. This lets an all-overflow batch backpropagate zero branch
gradients instead of failing because the output has no computational graph.
An expert that is never executed can have an unused gradient represented as
\texttt{None}. In comparisons, the tests interpret that as zero sensitivity,
not as proof that every optimizer treats an absent gradient identically to an
explicit zero tensor.

\subsection{Construct a deliberately expensive oracle}
\index{dense oracle}
A correctness oracle evaluates \emph{every} expert at every token, producing
$V\in\mathbb R^{T\times E\times D}$. Define
\begin{equation}
G_{tj}=a_t[\,k_t=j\,]p_{tj},\qquad
B_{td}=\sum_j G_{tj}V_{tjd}. \label{eq:denseoracle}
\end{equation}
This formula and sparse dispatch should agree on values and derivatives when
the routing plan is the same. The oracle wastes expert work but makes the
weighted sum easy to inspect.
\Listing{dense_oracle}

The tests compare outputs, input gradients, logit gradients, and all stacked
expert-parameter gradients at capacities $0,1,2,6$. A separate exact table tests
the plan itself. Sharing a plan between two execution paths does not verify
selection or admission, which is why both kinds of check are required.

An expert permutation provides another test. Away from ties, permute logit
columns and the corresponding expert weights together. The physical result
should stay the same. At ties, the first-index rule intentionally depends on
index order, so unrestricted permutation invariance is not promised.

\section{Differentiate the chosen computation, not a fictional smooth selector}
\index{routing gradient}\index{discrete decisions}
Within a region where $k_t$ and $a_t$ do not change, consider one admitted token
and write $b=p_k E_k(z)$. For upstream sensitivity
$\overline b=\partial\mathcal L/\partial b$, define the scalar
$s=\overline b^\top E_k(z)$. The softmax derivative is
\begin{equation}
\frac{\partial p_k}{\partial h_j}=p_k(\mathbf1_{j=k}-p_j).
\end{equation}
The chain rule then gives
\begin{equation}
\frac{\partial\mathcal L}{\partial h_j}
=a\,s\,p_k(\mathbf1_{j=k}-p_j). \label{eq:routergrad}
\end{equation}
Even an unselected expert's \emph{router logit} can receive a gradient because
it participates in the softmax denominator. That is different from evaluating
the unselected expert network or updating its weights through this token's
task loss.

For the expert parameters, the upstream sensitivity is scaled by $a p_k$.
For router weights, $\overline R=Z^\top\overline H$. For the input $z$, one
path passes through the selected expert and another through $h=zR$:
\begin{equation}
\overline z=a p_k J_{E_k}^{\top}\overline b
+R\,\overline h
\end{equation}
when writing individual feature sensitivities as column vectors. The dimensions
are $R:D\times E$, $\overline h:E$, and $\overline z:D$. In the complete block,
normalization and the identity residual add their own previously derived paths.

\Plate{03-gradient}{A routing boundary with no capacity pressure. For logits $(s,0)$ and constant expert outputs $2$ and $-1$, the selected output has smooth pieces but a jump at $s=0$. First-index ties choose expert 0, giving output 1 exactly at the boundary. The open circle marks the other branch's limiting value, $-0.5$.}{fig:gradient}

The argmax operation does not supply an ordinary derivative describing how a
small parameter change swaps experts. Away from a boundary it is locally
constant; at a tie the selected function may be discontinuous, as in
Figure \ref{fig:gradient}. Our autograd path differentiates the chosen smooth
piece and does not invent a straight-through estimator for the discrete
decision. Finite-difference checks use scores safely away from ties.

\subsection{Why top-1 renormalization can erase a learning path}
Suppose an implementation selects one expert and then renormalizes over only
the selected set. Its gate becomes
\begin{equation}
\widetilde p_k=\frac{p_k}{p_k}=1.
\end{equation}
The output is now $E_k(z)$. Away from selection boundaries, the task-loss
gradient through the router scores vanishes. This is not equivalent to
Equation \ref{eq:branch}. It might be a deliberate architecture with a separate
router objective, but it cannot silently replace probability weighting.

One test uses logits $(1,0,-1)$ and scalar expert value $2$. It checks
Equation \ref{eq:routergrad} directly, then verifies that the renormalized
expression has zero logit gradient. This negative control proves the test can
detect the missing router path. Gradient correctness is more informative here
than merely observing that the code runs.

\section{Capacity creates state and grouping semantics}
\index{capacity}\index{cohort}\index{chunk equivalence}
Let $d_j=\sum_t v_t[\,k_t=j\,]$ be demand in one cohort, with initially empty
slots. Under the stated policy,
\begin{equation}
q_j=\min(d_j,C),\quad
A=\sum_j q_j,\quad
D_{\mathrm{drop}}=T_v-A,\quad
U=EC-A, \label{eq:account}
\end{equation}
where $T_v=\sum_t v_t$ is the valid-token count. $q_j$ counts admitted tokens,
$A$ counts all admitted tokens, and $U$ counts unused slots. Demand is not
admitted load, and neither is total physical device work.

For $d=(3,2,1)$ and $C=2$, $A=5$, $D_{\mathrm{drop}}=1$, and $U=1$.
Increasing total capacity helps only when it changes a constrained expert's
available slots. Our capacity sweep is computed, not a hardware measurement.
\Plate{05-capacity}{Demand, overflow, and unused capacity are different quantities. At two slots per expert, one token is dropped despite one unused slot. Increasing capacity to three admits all six; further capacity adds unused slots under the same deterministic demand. Lines connect discrete integer-capacity cases.}{fig:capacity}
\begin{center}\input{\Generated/capacity-table.tex}\end{center}

A common family of designs sets capacity from a factor times average demand;
the Switch paper describes this trade-off \cite{switch}. For a concrete
integer policy one might choose $C=\lceil cT_v/E\rceil$. This reference instead
accepts $C$ explicitly so rounding and cohort size cannot remain hidden.
Changing $T_v$ or the capacity factor can change the forward computation.

\subsection{Resetting counters is not a harmless implementation detail}
Process the six-token example as a full cohort, then split after token 1. The
first two tokens both select expert 0 and fill it. If the second call starts
with fresh counters, token 2 is now admitted. If it receives the used counts
$(2,0,0)$, token 2 remains an overflow and the original result is recovered.

\Plate{04-cohort}{Capacity history under rechunking. Token logits and expert weights remain fixed. Resetting counters at the split changes admission of token 2; carrying used counts with the same capacity and order preserves the full-cohort decision. These counters are not an attention KV cache.}{fig:cohort}

The code copies incoming counters rather than mutating the caller's list.
Parity requires the same cohort, capacity, token order, logits, and parameters.
Carrying counters alone cannot repair logits changed by incorrect position or
attention handling. Nor does this demonstration establish an optimized
autoregressive cache implementation.

This reveals a boundary of Chapter 9's token-local FFN argument. Expert scoring
is token-local, but capacity admission couples tokens through a shared resource.
Reordering the first three tokens changes which one is dropped. Combining
otherwise isolated records into one capacity cohort can make one record's
admission depend on another record's demand, even if attention is correctly
record-masked.

With fixed capacity and first-arrival order, later tokens cannot change earlier
admissions. But computing capacity from a changing total length or selecting
admitted tokens by a global score ranking can introduce different dependencies.
A runtime must document grouping and admission order before claiming causal
or chunk-equivalent execution. If full-sequence training and token-at-a-time
inference use different capacity conventions, matching attention alone does
not prove model equivalence.

\section{Load balancing is an objective, not a certificate}
\index{load balancing}\index{auxiliary loss}
Use pre-admission demand fractions and mean router probabilities over valid
tokens:
\begin{equation}
f_j=\frac{1}{T_v}\sum_t v_t[\,k_t=j\,],\qquad
P_j=\frac{1}{T_v}\sum_t v_t p_{tj}.
\end{equation}
A Switch-style auxiliary term is
\begin{equation}
\mathcal L_{\mathrm{bal}}=\alpha E\sum_j f_jP_j,\qquad \alpha\ge0. \label{eq:balance}
\end{equation}
The assignment fractions $f_j$ are discrete and treated as fixed in the local
gradient. The probabilities $P_j$ are differentiable. Counting requests before
overflow matters: an overloaded expert should not look less demanded merely
because some requests were rejected.
\Listing{balance_loss}

Apply the softmax derivative and collect terms for a valid token:
\begin{align}
\frac{\partial\mathcal L_{\mathrm{bal}}}{\partial h_{tj}}
&=\frac{\alpha E}{T_v}\sum_r f_r p_{tr}
(\mathbf1_{r=j}-p_{tj})\\
&=\frac{\alpha E}{T_v}p_{tj}
\left(f_j-\sum_r f_rp_{tr}\right). \label{eq:balancegrad}
\end{align}
The tests compare this formula with autograd. An invalid token contributes no
term; a call with zero valid tokens raises an error rather than dividing by
zero. A dropped token can still receive an auxiliary router gradient even
though its expert branch contributed no task-loss gradient.

If $f_j=P_j=1/E$, Equation \ref{eq:balance} equals $\alpha$.
But this numeric value is not proof of uniform hard assignments. Set every
logit to zero. Then $P_j=1/E$, while first-index ties send every request to
expert 0, giving $f=(1,0,\ldots,0)$. The same loss value $\alpha$ appears.
Its gradient need not be zero. Therefore log actual demand, admissions, and
overflow along with the scalar loss. An attractive auxiliary-loss curve does
not certify balanced execution, expert specialization, or useful prediction.

\section{Replace the feed-forward branch in the existing block}
\index{pre-normalization}\index{residual stream}
The integration keeps Chapter 9's attention and normalizations:
\begin{align}
U&=X+\operatorname{Attn}(\operatorname{LN}_1(X)),\\
Z&=\operatorname{LN}_2(U),\\
Y_t&=U_t+a_t p_{t k_t}E_{k_t}(Z_t).
\end{align}
Our wrapper reuses \texttt{drafts/ch09/reference.py} for LayerNorm and attention
instead of reprinting a subtly different implementation. The review record
binds that reused source hash as well as the new chapter. The example fixes
two attention heads and disables dropout.
\Listing{sparse_block}

The full-compute oracle is also inserted at this same branch location. Tests
compare the integrated outputs and gradients through attention, normalization,
router, experts, and inputs. At zero capacity, the result must be exactly
$U$, not $X$ and not a zero tensor: the attention branch has already run.

There is also an exact dense limit. With one expert, its softmax probability
is one; if capacity admits all valid tokens and its weights equal the old FFN,
the unpadded no-dropout block equals Chapter 9's dense block. Increasing to
several experts changes the function and training problem; it does not merely
speed up the same already trained dense network.

\section{Count arithmetic without promising speed}
\index{multiply-accumulate count}\index{active parameters}
Ignoring bias additions and nonlinearities, one expert evaluation on one token
uses approximately $2DF$ multiply-accumulates (MACs). If $A$ tokens are admitted,
expert work is $2ADF$. Scoring all $T$ positions costs $TDE$ MACs in our
implementation, including padded positions whose scores are later ignored.
An all-expert oracle uses $2TEDF$ expert MACs.
\Listing{cost}

At $T=6,E=3,D=4,F=8,A=5$, the router uses $72$ MACs and admitted experts
use $320$, versus $1152$ expert MACs in the deliberately expensive oracle.
But the original \emph{single dense FFN} uses only $384$ MACs on six tokens.
The routed arithmetic total $392$ already exceeds that dense baseline, before
dispatch overhead. Confusing an all-expert oracle with the original dense
baseline would manufacture a misleading speedup story.

The experts store
\begin{equation}
E(2DF+F+D)=3(64+8+4)=228
\end{equation}
parameters, plus $DE=12$ router weights in this bias-free router. Stored
parameters, active arithmetic, optimizer state, and communication are separate
budgets. In a distributed design, gathering activations, transfers between
devices, imbalanced expert batches, synchronization, and padding can dominate
wall-clock time. None is measured by these MAC formulas.

Dropping work also changes the function. A low operation count obtained by
rejecting all tokens is not a successful optimization of a model that should
have evaluated their branches. Compare quality and semantics at matched
admission policies before measuring systems performance.

\section{What the reference establishes}
Run \texttt{python3 drafts/ch10/build.py --assets-only} to regenerate the
tables and \texttt{python3 -m pytest tests/test\_ch10\_reference.py} to run the
checks. The evidence covers a declared top-1 policy, literal routing examples,
forward/backward equivalence, deterministic ties, padding exclusion,
counter history, and a dense one-expert limit. It contains no genomic training
result or hardware benchmark.

Conditional expert routing is not itself biological regulation. In Evolutor's
taxonomy, Hermon is the Transformer-based direction and DOGMA the
non-Transformer direction; this chapter supplies a software mechanism that
must be specified and evaluated in any architecture that uses it. The upcoming
Chapters 1--10 integration pass will reconcile these execution contracts.
The next substantive biology-to-semantics chapter must explain which
biological observations motivate which abstractions, and where that mapping
fails.

\section{Exercises with worked reasoning}
\begin{enumerate}
\item \textbf{Count before running.} Demand is $(3,2,1)$ and $C=2$.
\emph{Solution:} admit $(2,2,1)$, drop one, leave one slot unused.
\item \textbf{Overflow output.} What remains for a dropped token?
\emph{Solution:} the incoming residual $u_t$, not a deleted token or necessarily the block's original $x_t$.
\item \textbf{Find the gate.} A chosen expert returns $(2,-1)$ and its probability is $0.7$.
\emph{Solution:} branch $(1.4,-0.7)$ before residual addition.
\item \textbf{Unselected logit.} Why can it receive a gradient?
\emph{Solution:} it changes the selected probability through the softmax denominator; it does not imply the unselected expert ran.
\item \textbf{Renormalization.} Differentiate $p_k/p_k$ within a fixed top-1 choice.
\emph{Solution:} the quotient is identically one, so its derivative is zero.
\item \textbf{Rechunking.} Split the example after token 1 and reset counters.
\emph{Solution:} token 2 changes from dropped to admitted. Carry $(2,0,0)$ with fixed capacity to preserve this policy.
\item \textbf{Tie permutation.} Can arbitrary expert relabeling preserve the result at an exact tie?
\emph{Solution:} not with index-based tie-breaking unless expert identity and tie order are also preserved.
\item \textbf{Load certificate.} Uniform probabilities give loss $\alpha$. Must hard demand be balanced?
\emph{Solution:} no; identical logits route every token to expert 0 under first-index ties.
\item \textbf{Dense limit.} State sufficient conditions for the original FFN result.
\emph{Solution:} one expert with matching weights, probability one, no overflow, matching valid inputs, attention and normalization, and no dropout.
\item \textbf{Arithmetic baseline.} Compare $72+320$ MACs with the original dense FFN's $384$.
\emph{Solution:} the routed count is eight larger; comparison with an all-expert oracle does not establish a dense-model speedup.
\item \textbf{Independent check.} Why is sparse/dense equality alone insufficient?
\emph{Solution:} both may share a wrong plan. Test explicit selection, demand, admitted IDs and capacity history independently.
\item \textbf{Evaluation design.} Does adding experts improve a genomic model?
\emph{Rubric:} specify held-out biological tasks, leakage controls, matched training and inference grouping, dropped-token rates, parameter and compute budgets, baselines, uncertainty, and measured hardware results. A working dispatch function answers none of these empirical questions by itself.
\end{enumerate}

\section*{Selective glossary}
\addcontentsline{toc}{section}{Selective glossary}
\textbf{Cohort:} tokens sharing one capacity budget and admission order.
\textbf{Demand:} requests before capacity rejection.
\textbf{Admission:} permission to execute the selected expert.
\textbf{Gate:} scalar weighting of an executed expert's output.
\textbf{Dispatch:} movement or indexing of token representations into and out of expert batches.
